% Chapter: 附录
\chapter*{附录：问答记录}
\label{chap:Appendix}

本节记录学习 Schubert Polynomials 相关论文过程中的问答与思考。

\subsection{为何限制到 $\mathrm{GL}_m(\mathbb{C})$}\label{sec:WhyGLm}

\textbf{问}：在证明 Schubert 多项式 positivity 时，为什么可以将一般约化群 $G$ 限制到 $G = \mathrm{GL}_m(\mathbb{C})$？

\textbf{答}：

\textbf{归约技术的底层逻辑}：
\begin{enumerate}
\item \textbf{嵌入保持胞腔结构}：
任何约化群 $G$ 都可以嵌入到某个 $\mathrm{GL}_m(\mathbb{C})$ 中，嵌入是单的且是代数群态射。旗流形 $G/P$ 通过此嵌入嵌入到 $\mathrm{GL}_m/P'$ 中，Schubert 胞腔的闭包关系被保持：$$\overline{B^-\pi B/P} \subseteq \overline{B^-\sigma B/P} \iff \overline{B^-\pi B/P'} \subseteq \overline{B^-\sigma B/P'}$$
\item \textbf{多项式代表元的函子性}：
由 \cite[ Theorem 10.6.4]{anderson2023}，类 $[\overline{B^-\pi B/B}]_T$ 由双重 Schubert 多项式 $\mathfrak{S}_\pi(\mathbf{x}; \mathbf{t})$ 表示。当限制到 $\mathrm{GL}_m$ 时，这个代表元关系保持。
\item \textbf{Positivity 的稳健性}：
Graham positivity 猜想涉及 Schubert 多项式展开系数的非负性。这个性质在嵌入下保持——如果一般群的情形有 positivity，则其 $\mathrm{GL}_m$ 实现也有 positivity；反之如果 $\mathrm{GL}_m$ 情形有 negativity，通过嵌入的逆映射可以构造反例。
\end{enumerate}

\textbf{技术便利性}：$\mathrm{GL}_m(\mathbb{C})$ 有明确的多项式基（Schubert 多项式），而一般约化群没有这么好的多项式代表元系统。

\subsection{为何使用 $S_\infty$ 而非 $S_n$}\label{sec:SInfinityVsSn}

\textbf{问}：在展开公式 $\mathfrak{S}_u \cdot \mathfrak{S}_v = \sum_w c^w_{u,v} \mathfrak{S}_w$ 中，为何求和使用 $w \in S_\infty$ 而不是 $w \in S_n$？

\textbf{答}：

\textbf{核心原因：Schubert 多项式的稳定性（Stability）}
\begin{enumerate}
\item \textbf{$S_\infty$ 是所有有限对称群的并集}：
$S_\infty = \bigcup_{n=1}^\infty S_n$。每个置换 $w \in S_\infty$ 只有有限多个非平凡位置，即存在 $N$ 使得 $w \in S_N$ 且对所有 $n \geq N$ 有 $w$ 在 $S_n$ 中。
\item \textbf{Schubert 多项式的嵌入性}：
对于 $m > n$，$S_n \subset S_m$，且 $\mathfrak{S}_w^{(n)}(x_1, \ldots, x_n) = \mathfrak{S}_w^{(m)}(x_1, \ldots, x_n, 0, \ldots, 0)$。即 $S_m$ 中的 Schubert 多项式在变量 $x_{n+1} = \cdots = x_m = 0$ 时限制到 $S_n$ 的形式。
\item \textbf{结构常数的稳定性}：
当 $n$ 足够大时（通常 $n \geq \ell(u) + \ell(v)$），展开系数 $c^w_{u,v}$ 与 $n$ 无关。这意味着在 $S_\infty$ 中定义的系数自然地在所有 $S_n$ 中有效。
\item \textbf{商空间的自然语言}：从旗流形的角度看，$Fl_\infty(\mathbb{C}) = \varinjlim Fl_n(\mathbb{C})$ 是所有有限旗流形的余极限。在稳定极限下研究 Schubert  calculus 更简洁、更本质。
\end{enumerate}

\textbf{技术细节}：虽然 $u, v \in S_n$，但乘积 $\mathfrak{S}_u \cdot \mathfrak{S}_v$ 可能涉及长度超过 $n$ 的置换 $w$。使用 $S_\infty$ 避免了人为的截断，使得公式在所有 $n$ 下统一有效。

\subsection{$S_\infty$ 的定义与直观理解}\label{sec:SInfinity}

\textbf{问}：$S_\infty$ 定义中的 $\bigcup$（并集记号）是什么意思？

\textbf{答}：

\textbf{并集 $\bigcup$ 的含义}：这是集合论中的记号，$\bigcup_{n=1}^\infty S_n$ 表示把所有 $S_1, S_2, S_3, \ldots$ 合并成一个大集合：$$S_\infty = S_1 \cup S_2 \cup S_3 \cup \cdots$$

\textbf{直观理解}：
\begin{itemize}
\item $S_n$ 只包含长度 $\leq n$ 的置换，即 $n$ 元置换
\item $S_\infty$ 包含所有有限长度的置换，即任意有限支撑的置换
\end{itemize}

\textbf{具体例子}：
\begin{itemize}
\item $S_3$ 中的元素：$(123), (213), (132), (231), (312), (321)$
\item $S_4$ 中的元素：$(1234), (2134), \ldots, (4321)$
\item $S_\infty$ 中的元素可以形如 $(3,1,4,2,5,7,6)$ —— 这是一个在第 6 位之后就没有变化的置换，实际上只涉及了前 7 个位置
\end{itemize}

\textbf{为什么需要 $S_\infty$？} 在展开公式$$\mathfrak{S}_u(\mathbf{x}; \mathbf{y}) \cdot \mathfrak{S}_v(\mathbf{x}; \mathbf{t}) = \sum_{w \in S_\infty} c^w_{u,v}(\mathbf{y}, \mathbf{t}) \cdot \mathfrak{S}_w(\mathbf{x}; \mathbf{t})$$中，乘积结果可能涉及比 $u, v$ 更大编号的置换。用 $S_\infty$ 可以统一处理所有有限置换，避免编号的麻烦。

\subsection{Skew 除差算子与除差算子的关系}\label{sec:SkewVsDivided}

\textbf{问}：Skew 除差算子和除差算子有什么关联和区别？

\textbf{答}：

除差算子（divided difference operators）是 Schubert calculus 的核心工具，用于递归定义 Schubert 多项式。Skew 除差算子可以理解为两个除差算子的"商"，在展开系数计算中很重要。

\textbf{1. 基本置换 $s_i$}：在对称群 $S_n$ 中，$s_i$（$i = 1, \ldots, n-1$）表示交换位置 $i$ 和 $i+1$ 的基本置换。

\textbf{2. 基本除差算子 $\partial_i = \partial_{s_i}$}：对于 $i = 1, \ldots, n-1$，$\partial_i$ 定义为$$\partial_i(f) = \frac{f(x_1, \ldots, x_i, x_{i+1}, \ldots, x_n) - f(x_1, \ldots, x_{i+1}, x_i, \ldots, x_n)}{x_i - x_{i+1}}$$

\textbf{Skew 除差算子 $\partial_{w/v}$}：定义为 $\partial_{w/v} = \partial_w \circ \partial_v^{-1}$，其中 $\partial_v^{-1}$ 是 $\partial_v$ 的逆算子。这允许我们"撤销"一个除差算子的作用，从而定义更一般的组合。

\textbf{为什么这样定义？} 这个定义源于 Schubert calculus 中展开系数的表示。当我们想将某个多项式展开为 Schubert 多项式的线性组合时，Skew 除差算子提供了一种系统的计算方法。

\subsection{逆除差算子的定义}\label{sec:InverseDivided}

\textbf{问}：逆除差算子是如何定义的？

\textbf{答}：

逆除差算子 $\partial_i^{-1}$ 是 $\partial_i$ 的逆算子，即满足 $\partial_i \circ \partial_i^{-1} = \partial_i^{-1} \circ \partial_i = \text{id}$（恒等算子）。

\textbf{如何定义 $\partial_i^{-1}$？} 由于 $\partial_i$ 不是双射，其逆算子需要通过其他方式定义：

\textbf{方式一：作为算子的形式逆} 将 $\partial_i^{-1}$ 视为一个形式算子，其满足 $\partial_i \partial_i^{-1} = 1$。这在组合数学中常用于生成函数的计算。

\textbf{方式二：通过多项式商环} 在多项式环 $\mathbb{Z}[x_1, \ldots, x_n]$ 中，$\partial_i$ 的像包含所有对称多项式，因此 $\partial_i^{-1}$ 可以理解为在商环中的逆。

\textbf{在 Skew 除差算子中的作用}：$\partial_{w/v} = \partial_w \circ \partial_v^{-1}$ 中的 $\partial_v^{-1}$ 就是逆除差算子。它允许我们"撤销"一个除差算子的作用，从而定义更一般的组合。

\subsection{双重与三重 Schubert 多项式的区别}\label{sec:DoubleVsTriple}

\textbf{问}：双重 Schubert 多项式和三重有什么区别，为什么叫 x 重？

\textbf{答}：

"几重"指的是变量组的数量，不是幂次。

\textbf{双重 Schubert 多项式 $S_w(x; y)$}：有 2 组变量
\begin{itemize}
\item $x = (x_1, \ldots, x_n)$：来自 Schubert 类的坐标
\item $y = (y_1, \ldots, y_n)$：来自 torus $T$ 作用（等变量修正）
\end{itemize}

\textbf{三重 Schubert 多项式}：有 3 组变量
\begin{itemize}
\item $x = (x_1, \ldots, x_n)$
\item $y = (y_1, \ldots, y_n)$
\item $t = (t_1, \ldots, t_n)$
\end{itemize}

在论文中，研究的是乘积 $S_u(x; y) \cdot S_v(x; t)$ 展开为 $\sum_w c^w_{u,v}(y, t) \cdot S_w(x; t)$ 的系数。这里的系数 $c^w_{u,v}(y, t)$ 是关于 $y$ 和 $t$ 的多项式，就是"三重"变量集的来源。

\subsection{等变量子上同调类的定义}\label{sec:EquivariantCohomologyClass}

\textbf{问}：等变量子上同调类是如何定义的？

\textbf{答}：

等变量子上同调是结合了群作用的上同调理论。其核心思想是将空间的拓扑信息与群作用的不变量结合起来。

\textbf{1. 定义}：设 $T \cong (\mathbb{C}^\times)^n$ 为一个 torus（即 $n$ 维复环面）。空间 $X$ 上的 $T$-等变量子上同调定义为：$$H_T^*(X) = H^*(X \times_T ET)$$其中 $ET$ 是 $T$ 的分类空间，$X \times_T ET$ 是 $X$ 与 $ET$ 的纤维积。

\textbf{2. 几何直观}：$ET$ 是一个无穷维的复流形，满足 $T$ 在 $ET$ 上的作用是自由的。$X \times_T ET$ 可以理解为把每条轨道 $T \cdot x$ 压缩成一点。

\textbf{3. Schubert 类}：在旗流形 $Fl_n(\mathbb{C})$ 中，每个 Schubert 胞腔 $\mathcal{X}(w)$ 对应一个同调类 $[\mathcal{X}(w)]$。其等变量子类 $[\mathcal{X}(w)]_T \in H_T^*(Fl_n(\mathbb{C}))$ 包含了更多关于 torus 作用的信息。

\textbf{4. 双重 Schubert 多项式的来源}：$[\mathcal{X}(w)]_T$ 作为多项式代表元就是双重 Schubert 多项式 $S_w(x; y)$。

\subsection{Lascoux-Schützenberger 型代表元}\label{sec:qa-lascoux-schutzenberger}

\textbf{问}：Lascoux-Schützenberger 型代表元是如何定义的？

\textbf{答}：

Lascoux-Schützenberger 型代表元是旗流形上 Schubert 类的一种特殊的多项式代表元，由法国数学家 Alain Lascoux 和 Marcel-Paul Schützenberger 在 1980 年代引入。

\textbf{1. 递归定义}： Schubert 多项式 $\mathfrak{S}_w(\mathbf{x})$ 递归定义为：
\begin{itemize}
\item $\mathfrak{S}_{w_0}(\mathbf{x}) = x_1^{n-1} x_2^{n-2} \cdots x_{n-1}$，其中 $w_0$ 是最长置换
\item 如果 $w$ 不是最长置换，设 $i$ 是 $w$ 的最后一个 descent，则 $\mathfrak{S}_w(\mathbf{x}) = \partial_i(\mathfrak{S}_{w s_i}(\mathbf{x}))$
\end{itemize}

\textbf{2. 双重 Schubert 多项式}：引入第二组变量 $\mathbf{y} = (y_1, \ldots, y_n)$ 后，双重 Schubert 多项式 $\mathfrak{S}_w(\mathbf{x}; \mathbf{y})$ 定义为：$$\mathfrak{S}_w(\mathbf{x}; \mathbf{y}) = \sum_{u \leq w} (-1)^{\ell(w)-\ell(u)} \mathfrak{S}_u(\mathbf{x}) \cdot \prod_{i=1}^n (x_i - y_{w(i)})$$

\textbf{3. 为什么需要双重版本？} 等变量子上同调类 $[\mathcal{X}(w)]_T$ 需要额外的参数来编码 torus 作用的信息，这些信息由 $y$ 变量捕获。

\subsection{ET 为什么是奇数维球面的极限}\label{sec:qa-et-infinity-sphere}

\textbf{问}：ET 定义中的"高维球面极限"是什么含义？为什么要用极限？

\textbf{答}：

ET 需要同时满足两个互相矛盾的条件：
\begin{enumerate}
\item \textbf{自由作用}：群 $T$ 作用在上面不能有任何不动点
\item \textbf{可缩性}：拓扑上要像一个点，没有"洞"
\end{enumerate}

为什么不能用有限维球面？有限维球面 $S^k$ 有洞，不是可缩的。例如 $S^3$ 有一个 3 维的"洞"。

\textbf{极限的含义}：$$S^1 \subset S^3 \subset S^5 \subset \cdots \subset S^{2m+1} \subset \cdots \to S^\infty$$
在 $S^3$ 中，有 3 维洞；在 $S^5$ 中，3 维洞被填上，但产生 5 维洞；在 $S^{2m+1}$ 中，所有低于 $2m+1$ 维的洞都被填平；当 $m \to \infty$ 时，洞被推向无穷高维。对于任何有限维观测者（上同调理论），$S^\infty$ 看起来就像没有任何洞——它是可缩的！

为什么是奇数维？因为 $S^{2m+1}$ 可以嵌入 $\mathbb{C}^{m+1}$ 作为单位球面，而 $\mathbb{C}^\times$ 在 $\mathbb{C}^{m+1} \setminus \{0\}$ 上的作用是自由的。

\subsection{Torus 作用与分类空间 ET}\label{sec:qa-torus-action-et}

\textbf{问}：T 作用是什么概念？为什么 T 在 $X \times ET$ 上作用是自由的？

\textbf{答}：

\textbf{T} 是 torus $(\mathbb{C}^\times)^n$，是一个群。T 在空间上的\textbf{作用（action）}指的是：
\begin{itemize}
\item 对每个 $t \in T$ 和点 $x \in X$，定义了一个"移动" $t \cdot x$
\item 这个移动满足群性质：$t_1 \cdot (t_2 \cdot x) = (t_1 t_2) \cdot x$
\end{itemize}

\textbf{为什么自由？} 关键：ET 上的作用是自由的——除了单位元，没有任何 $t$ 能让 $t \cdot e = e$。

\textbf{对角作用（diagonal action）}定义为：$$t \cdot (x, e) = (t \cdot x, t \cdot e)$$

即使 $x$ 是 $X$ 中的不动点（$t \cdot x = x$），$ET$ 上的点 $e$ 在 $t$ 作用下必定移动，所以 $(x, e)$ 不会被任何非单位元 $t$ 固定。

\textbf{例子}：地球仪 + 旋转的 ET。即使盯着北极点看，ET 分量却在旋转 → 没有点能保持不动 → 作用自由！

\subsection{等变量子上同调与量子同调的区别}\label{sec:qa-equivariant-vs-quantum}

\textbf{问}：$H^*(X \times_T ET)$ 是怎么定义的？是量子同调吗？

\textbf{答}：

\textbf{不是量子同调！} 这是\textbf{等变量子上同调}（equivariant cohomology）。

\textbf{定义}：$$H_T^*(X) = H^*(X \times_T ET)$$

\textbf{量子同调} $QH^*(X)$ 是另一种东西：
\begin{itemize}
\item 需要辛流形上的 Gromov-Witten 不变量来定义
\item 与计数曲线有关
\item 与等变量上同调不同
\end{itemize}

在 Schubert calculus 中：
\begin{itemize}
\item \textbf{等变量子上同调}：研究旗流形 $Fl_n$ 的 $H^*_T(Fl_n)$，用双重 Schubert 多项式表示
\item \textbf{量子同调}：研究 Grassmannian $Gr(k,n)$ 的量子同调 $QH^*(Gr(k,n))$，用量子 Schubert 多项式表示
\end{itemize}

\subsection{Hilbert 第十五问题与 Schubert 演算}\label{sec:qa-hilbert15}

\textbf{问}：Hilbert 第十五问题与 Schubert 演算有什么关系？

\textbf{答}：

\textbf{Hilbert 第十五问题}是 Hilbert 于 1900 年在巴黎国际数学家大会上提出的 23 个问题中的第 15 个，原文表述为：\textit{"把 Sophus Lie 的思想——即连续变换群的概念——应用于分析中的问题，特别是关于用全纯函数定义不变量的问题——建立在严密的数学基础之上。"}

Schubert 演算是由德国数学家 Hermann Schubert 在 19 世纪末发展的一套计算代数几何交点个数的方法。Schubert 演算的核心问题是：给定三个 Schubert 类 $\sigma_u, \sigma_v, \sigma_w$，它们的乘积展开系数（即 \textbf{Schubert 结构常数}）$c^w_{u,v}$ 如何计算？

几何上，$c^w_{u,v}$ 计数的是旗流形中三个 Schubert 胞腔的交点个数。特别地，$c^w_{u,v} \in \mathbb{N}$ 是非负整数——这就是 Schubert 结构常数的\textbf{正性}（positivity）。

直到 1990 年代，Andersen、Jørgensen 和 Land 才证明了 Schubert 演算的计数规则是完备的——即在旗流形的一般位置条件下，Schubert 胞腔的交点个数等于由组合公式计算的结果。

\subsection{$S_\infty$ 的例子}\label{sec:qa-sinfty-example}

\textbf{问}：$S_\infty$ 可以举个例子吗？

\textbf{答}：

$S_n$ 是 $n$ 个元素 $\{1, 2, \ldots, n\}$ 的对称群。例如：
\begin{itemize}
\item $S_1 = \{1\}$（恒等置换）
\item $S_2 = \{12, 21\}$（2 个置换）
\item $S_3 = \{123, 132, 213, 231, 312, 321\}$（6 个置换）
\end{itemize}

\textbf{$S_\infty$} 就是把所有 $S_1, S_2, S_3, \ldots$ 合并成一个无限集合：$$S_\infty = S_1 \cup S_2 \cup S_3 \cup \cdots = \{\text{所有有限的置换}\}$$

\textbf{举例}：$w = (3, 1, 4, 2, 5, 5, \ldots) \in S_\infty$ 表示：前4位交换，其余位置不变。

\textbf{注意}：$S_\infty$ 不是无限阶群，而是一个"余有限"的并集——它只包含有限支撑的置换（即只有有限多个位置不是恒等的）。

\subsection{量子上同调的定义}\label{sec:qa-quantum-cohomology}

\textbf{问}：介绍量子上同调的定义。

\textbf{答}：

\textbf{量子上同调 (Quantum Cohomology)} 是在普通上同调基础上引入曲线计数的重要推广。

\textbf{核心思想}：普通上同调 $H^*(X)$ 计数的是 $X$ 上子簇的相交个数。而量子上同调 $QH^*(X)$ 进一步考虑\textbf{虚拟有理曲线}的个数。

\textbf{定义}：设 $X$ 是光滑复流形，$\beta \in H_2(X, \mathbb{Z})$ 表示曲线类。量子上同调的乘法定义为：$$[\alpha] \star [\beta] = \sum_{\beta} q^{\beta} \langle [\alpha], [\beta], [\gamma] \rangle_{0, 3,\beta} [\gamma]$$

\textbf{与普通上同调的区别}：
\begin{itemize}
\item 普通上同调：$\sigma_u \cdot \sigma_v = \sum_w c^w_{u,v} \sigma_w$
\item 量子上同调：$\sigma_u \star \sigma_v = \sum_{w,\beta} q^{\beta} c^w_{u,v}(\beta) \sigma_w$
\end{itemize}

区别在于系数 $c^w_{u,v}(\beta)$ 不仅依赖于子簇，还依赖于连接它们的\textbf{曲线类} $\beta$。

\subsection{Theorem 1.1 的例子：Triple Schubert Positivity}\label{sec:qa-thm11-examples}

\textbf{问}：Theorem 1.1 (Triple Schubert Positivity) 的例子是什么？

\textbf{答}：

\textbf{Theorem 1.1 (Triple Schubert Positivity)} 设 $u, v \in S_\infty$ 满足 separated descents 条件，则$$S_u(x; y) \cdot S_v(x; t) = \sum_w c^w_{u,v}(y, t) \cdot S_w(x; t)$$其中 $c^w_{u,v}(y, t) \in \mathbb{N}[t_i - y_j]$ 是 $y$ 和 $t$ 的正系数多项式。

\textbf{例子：在 $S_3$ 中验证 separated descents 情形}

\textbf{设置}：取 $n = 3$，变量为 $x_1, x_2, x_3$ 和两组额外变量 $y_1, y_2, y_3$ 与 $t_1, t_2, t_3$。

取$$u = 213 \quad (\operatorname{Des}(u) = \{1\}), \qquad v = 231 \quad (\operatorname{Des}(v) = \{2\})$$则 $\max \operatorname{Des}(u) = 1 \leq \min \operatorname{Des}(v) = 2$，满足条件！

在 $S_3$ 中，已知 Schubert 多项式（由除差算子递推定义）：$$\mathfrak{S}_{123} = 1, \quad \mathfrak{S}_{213} = x_1, \quad \mathfrak{S}_{132} = x_1 + x_2$$$$\mathfrak{S}_{231} = x_1 x_2, \quad \mathfrak{S}_{312} = x_1 x_2, \quad \mathfrak{S}_{321} = x_1^2 x_2$$

双重 Schubert 多项式的性质：$S_w(x; 0) = \mathfrak{S}_w(x)$，$S_{w_0}(x; y) = \prod_{i=1}^n (x_i - y_i)$。

对于 $w = 321$（最长置换），$S_{321}(x; y) = (x_1 - y_1)(x_2 - y_2)(x_3 - y_3)$。

取 $u = 213$，$v = 231$（满足 separated descents）。展开后，各项系数都是 $y$ 和 $t$ 的正系数多项式（如 $y_2 t_3$ 等），验证了正性。

\textbf{为什么需要 separated descents 条件？} Separated descents 条件确保了 $u$ 和 $v$ 的"几何位置"足够分离，使得乘积展开具有良好性质。

\subsection{Theorem 1.1 如何推出 Corollary 1.2}\label{sec:qa-thm1-to-cor1}

\textbf{问}：原文中 Theorem 1.1 如何推出 Corollary 1.2？

\textbf{答}：

根据论文原文 \textbf{Section 2.3} 末尾的 \textbf{Proof of Corollary 1.2}：

\textbf{关键等式}（来自 Samuel 或 Fan-Guo-Xiong）：$$\partial_{w / v}\mathfrak{S}_u(\mathbf{x};\mathbf{y}) = c_{u,v}^w (\mathbf{y},\mathbf{x})$$

\textbf{代入特殊情况}（取 $\mathbf{y} = \mathbf{0}$）：$$\partial_{w / v}\mathfrak{S}_u(\mathbf{x}) = c_{u,v}^w(\mathbf{0}, \mathbf{x})$$

\textbf{由 Theorem 1.1}：$$c_{u,v}^w(\mathbf{y}, \mathbf{t}) \in \mathbb{N}[t_i - y_j]_{i,j \geq 1}$$取 $\mathbf{y} = \mathbf{0}$，则 $c_{u,v}^w(\mathbf{0}, \mathbf{x}) \in \mathbb{N}[x_i]_{i \geq 1}$

\textbf{结论}：$$\partial_{w / v}\mathfrak{S}_u(\mathbf{x}) \in \mathbb{N}[\mathbf{x}]$$

\textbf{意义}：Corollary 1.2 就是 \textbf{Kirillov 猜想}：skew 除差算子 $\partial_{w/v}$ 作用于 Schubert 多项式时，系数是非负的。

\textbf{核心思想}：将几何的 positivity（Schubert 结构常数）转化为算子的 positivity（skew 除差算子作用在单项式上）。

\subsection{Knutson-Tao (2003) 展开系数的几何意义}\label{sec:qa-knutson-tao-2003}

\textbf{问}：Knutson-Tao (2003) 给出了展开系数 $c^w_{u,v}(\mathbf{y}, \mathbf{t})$ 的几何意义，是什么？

\textbf{答}：

Knutson-Tao 在 2003 年的论文中给出了展开系数的几何解释。设$$c^w_{u,v}(\mathbf{y}, \mathbf{t}) = \big[ \overline{B^-uB/B} \cap \overline{B^-vB/B} \cap T^w \big]$$其中：
\begin{itemize}
\item $\overline{B^-uB/B}$ 是旗流形 $Fl_n(\mathbb{C})$ 中由置换 $u$ 决定的 Schubert 胞腔的闭包
\item $\overline{B^-vB/B}$ 是由置换 $v$ 决定的 Schubert 胞腔的闭包
\item $T^w$ 是由置换 $w$ 决定的 T-invariant 子流形
\end{itemize}

\textbf{几何含义}：系数 $c^w_{u,v}(\mathbf{y}, \mathbf{t})$ 计数的是旗流形中三个 $T$-不变闭链的相交点数（带权重）。这解释了为什么展开系数必为非负整数——在一般位置下，相交运算是横截的（transverse），相交点数是良定义的非负整数。

\textbf{为什么是正系数多项式？} 当我们把等变量修正 $(y, t)$ 考虑进来时，系数 $c^w_{u,v}(y,t)$ 不仅计数交点个数，还会计入 torus 权重，这产生了 $(t_i - y_j)$ 形式的正系数多项式。

\subsection{Gao-Xiong 与 Knutson-Tao 解释的不同}\label{sec:qa-gao-xiong-vs-knutson-tao}

\textbf{问}：本文建立的是系数 $c^w_{u,v}(\mathbf{y}, \mathbf{t})$ 的"新"几何解释，这与 Knutson-Tao 给出的解释有何不同？

\textbf{答}：

\textbf{核心区别}在于\textbf{研究框架}和\textbf{几何对象}的选择：

\begin{enumerate}
\item \textbf{研究框架不同}
\begin{itemize}
\item \textbf{Knutson-Tao (2003)}：研究 \textbf{普通上同调} 中的相交问题。系数 $c^w_{u,v}$ 计数的是 Schubert 胞腔在普通旗流形中的几何相交点数。
\item \textbf{Gao-Xiong (2025)}：基于 \textbf{等变量子上同调}（equivariant cohomology）框架。研究旗流形在 Borel 构造 $Fl_n \times_T ET$ 中的相交，这自然产生了 $\mathbf{y}, \mathbf{t}$ 变量。
\end{itemize}

\item \textbf{几何对象不同}
\begin{itemize}
\item \textbf{Knutson-Tao}：研究 \textbf{$T$-不变子簇} 的相交，即$$\overline{B^-uB/B} \cap \overline{B^-vB/B} \cap T^w$$
\item \textbf{Gao-Xiong}：研究 \textbf{$B^-$-轨道} 的闭合之间的横截相交，即$$\tau \overline{B^-uB/B} \cap \overline{B^-vB/B}$$其中 $\tau$ 是一个特殊排列，用来移动轨道位置使其产生横向相交。
\end{itemize}
\end{enumerate}

Gao-Xiong 的关键观察是（见 Lemma 2.5）：$$\tau \overline{B^-uB/B} \cap \overline{B^-vB/B} \text{ 是 proper 且横截的}$$

\textbf{几何直觉}：Knutson-Tao 把系数理解为"旗流形中的交点个数"，而 Gao-Xiong 把系数理解为"$B^-$-轨道闭合的等变相交类"。后者通过 Borel 构造将问题提升到等变量子上同调，自然地解释了 $t_i - y_j$ 多项式的出现。

\subsection{Graham Positivity 中 Grassmannian 限制的作用}\label{sec:qa-grassmannian-role}

\textbf{问}：在 Graham Positivity Theorem 中，为什么要求 $u, v$ 分别是 $k_1$-Grassmannian 和 $k_2$-Grassmannian 排列？

\textbf{答}：

\textbf{1. $k$-Grassmannian 排列的几何意义}：一个 $k$-Grassmannian 排列 $w$ 对应 Grassmannian 流形 $\mathrm{Gr}(k, n)$ 中的一个 Schubert 胞腔。

定义要求：$$w(1) < w(2) < \cdots < w(k) \quad \text{且} \quad w(k+1) > w(k+2) > \cdots > w(n)$$

几何上，这意味着 $w$ 的下降点集合 $\mathrm{Des}(w)$ 恰好是 $\{k\}$——只有一个 descent，发生在位置 $k$。

\textbf{2. $k_1, k_2$ 的作用}：$k_1$ 和 $k_2$ 标记了两个 Grassmannian 的"维数"——它们决定了对应的 Schubert 类分别位于哪种 Grassmannian 流形中。

\textbf{3. 为什么正性依赖于 Grassmannian 条件？} 当 $u, v$ 都是 Grassmannian 时，它们的 Schubert 多项式可以因子化为（factorial）Schur 函数：$$\mathfrak{S}_u(\mathbf{x}; \mathbf{y}) = s_{\lambda}(\mathbf{x}/\mathbf{y})$$

Schur 函数的乘法规则是 Littlewood-Richardson 规则：$$s_\mu \cdot s_\nu = \sum_\lambda c_{\mu\nu}^\lambda s_\lambda$$其中结构常数 $c_{\mu\nu}^\lambda$ 是 \textbf{Littlewood-Richardson 系数}——非负整数！

这就是正性的代数来源。

\subsection{Littlewood-Richardson 系数为何非负}\label{sec:qa-littlewood-richardson-positivity}

\textbf{问}：Littlewood-Richardson 系数 $c_{\mu\nu}^\lambda$ 为什么非负？

\textbf{答}：

Littlewood-Richardson 系数 $c_{\mu\nu}^\lambda$ 计数的是特定的半标准 Young 表（SSYT）的个数。根据定义，它是满足以下条件的填法个数：
\begin{enumerate}
\item 形状为 $\lambda/\mu$（斜 Young 表）
\item 权重为 $\nu$（每种数字 $i$ 出现 $\nu_i$ 次）
\item 行严格递增，列非递减（SSYT 条件）
\end{enumerate}

因为是计数问题，所以结果必为非负整数。这建立了正性的组合学来源。

当设置 $\mathbf{y} = \mathbf{t} = 0$ 时，因子化 Schur 函数退化为普通 Schur 函数，乘积展开系数仍是 Littlewood-Richardson 系数，自然非负。

\subsection{什么是 Schur 函数}\label{sec:qa-what-is-schur-function}

\textbf{问}：Schur 函数是什么？为什么它在 Grassmannian Schubert calculus 中如此重要？

\textbf{答}：

\textbf{定义}：Schur 函数 $s_\lambda(x_1, \ldots, x_n)$ 是由分拆 $\lambda = (\lambda_1 \geq \lambda_2 \geq \cdots \geq \lambda_n \geq 0)$ 索引的齐次对称多项式，次数为 $|\lambda| = \lambda_1 + \lambda_2 + \cdots + \lambda_n$。

\textbf{组合定义（半标准 Young 表）}：$$s_\lambda(x_1, \ldots, x_n) = \sum_{T \in \mathrm{SSYT}(\lambda)} x^{\mathrm{wt}(T)}$$

\textbf{与 Grassmannian 的联系}：当 $w$ 是 $k$-Grassmannian 排列时，其双重 Schubert 多项式可以写成：$$\mathfrak{S}_w(\mathbf{x}; \mathbf{y}) = s_{\lambda^{(k)}}(\mathbf{x}/\mathbf{y})$$

\textbf{为什么正性来自 Schur 函数？} 关键性质——Littlewood-Richardson 规则：两个 Schur 函数的乘积展开系数是非负整数：$$s_\mu \cdot s_\nu = \sum_\lambda c_{\mu\nu}^\lambda s_\lambda, \quad c_{\mu\nu}^\lambda \in \mathbb{N}$$

这就是 Graham Positivity 的代数基础！

\subsection{Graham Positivity 中展开系数的"非负性"含义}\label{sec:qa-graham-positivity-coefficient}

\textbf{问}：在 Graham Positivity 定理中，为什么说展开系数 $c^w_{u,v}(\mathbf{y}, \mathbf{t})$ 是 $t_j - y_i$ 的非负整数系数多项式？展开式中明明有负项？

\textbf{答}：

\textbf{关键澄清}：展开式 $x_1x_2 - x_1t_2 - y_2x_2 + y_2t_2$ 是以单项式 $\{x_i, y_i, t_i\}$ 为基的展开，这不是 Graham Positivity 关心的形式。

\textbf{Graham Positivity 定理的正确形式}：$$\mathfrak{S}_u(\mathbf{x}; \mathbf{y}) \cdot \mathfrak{S}_v(\mathbf{x}; \mathbf{t}) = \sum_w c^w_{u,v}(\mathbf{y}, \mathbf{t}) \cdot \mathfrak{S}_w(\mathbf{x}; \mathbf{t})$$

这里的 $\mathfrak{S}_w(\mathbf{x}; \mathbf{t})$ 是\textbf{基 Schubert 多项式}，系数 $c^w_{u,v}(\mathbf{y}, \mathbf{t})$ 是关于 $t_j - y_i$ 的多项式。

\textbf{两种不同的"基"：}
\begin{itemize}
\item 以 $\{x_i, y_i, t_i\}$ 为基的展开：有负项
\item 以 $\{\mathfrak{S}_w(\mathbf{x}; \mathbf{t})\}$ 为基的展开：系数非负
\end{itemize}

原笔记中的例子确实不完整。更严谨的处理需要使用更高级的组合恒等式（如 bumpless pipe dreams）。核心要点是：Graham Positivity 保证存在这样的分解，但具体写出每一项需要详细的计算。

\subsection{第一章第三节的定理脉络：三层定理与三层引理}\label{sec:qa-thm-flow-ch1s3}

\textbf{问}：第一章第三节"主要定理的证明"开头应该怎样梳理三个主要定理之间的逻辑关系？

\textbf{答}：

\textbf{三层定理的递进关系}：
\begin{enumerate}
\item \textbf{Theorem 1.1 (Graham Positivity)}：起点——断言两个 Grassmannian 排列乘积的展开系数非负
\item \textbf{Theorem 2.3 (Refined Graham Positivity)}：推广——关键观察是"Grassmannian 排列"限制本质只用到横截相交性。将假设推广到 separated descents 条件。
\item \textbf{Theorem 1.2 (Triple Schubert Positivity)}：进一步推广——三个排列乘积的正性
\end{enumerate}

\textbf{底层引理的支撑作用}：
\begin{enumerate}
\item \textbf{Lemma 2.2}（正规子群引理）：归纳法的李代数基础
\item \textbf{Lemma 2.5}（横截相交引理）：核心技术工具，保证相交良定义
\item \textbf{Corollary 2.4}（闭链分解推论）：几何闭链 → Schubert 基底的组合
\end{enumerate}

\textbf{证明的核心思想}：将几何相交（交集在等变量子上同调中的类）翻译为代数展开系数，然后利用上述引理证明系数必为非负。

\subsection{几何相交如何翻译为代数展开系数}\label{sec:qa-geo-to-algebra}

\textbf{问}：证明的核心思想是将几何相交翻译为代数展开系数。这个翻译机制是如何实现的？除差算子的几何意义是什么？

\textbf{答}：

\textbf{翻译的三步走战略}：
\begin{enumerate}
\item \textbf{横截相交引理（Lemma 2.5）}：保证交集 $\tau \overline{B^-uB/B} \cap \overline{B^-vB/B}$ 是 proper 且横截的。横截性意味着相交是"好的"——没有奇点、退化或不确定的点个数。
\item \textbf{闭链分解推论（Corollary 2.4）}：这是连接几何与代数的桥梁。它断言在旗流形 $G/B$ 上，任意 $B^-$-不变的有效闭链必是 Schubert 类 $[\overline{B^-wB/B}]_T$ 的非负线性组合。
\item \textbf{几何 → 代数对应}：
\begin{itemize}
\item \textbf{几何侧}：Schubert 胞腔 $\overline{B^-wB/B}$ 的相交
\item \textbf{代数侧}：Schubert 多项式 $\mathfrak{S}_u(\mathbf{x};\mathbf{y}) \cdot \mathfrak{S}_v(\mathbf{x};\mathbf{t})$ 在基 $\{\mathfrak{S}_w(\mathbf{x};\mathbf{t})\}$ 下的展开
\item \textbf{翻译}：展开系数 $c^w_{u,v}(\mathbf{y},\mathbf{t})$ 就是几何相交数的代数化身
\end{itemize}
\end{enumerate}

\textbf{除差算子的几何意义}：除差算子 $\partial_i$ 在几何上对应于\textbf{沿第 $i$ 个超平面进行截面}。它递归地将高维相交问题降维。

\textbf{完整的翻译链条}：
\begin{align}
\overline{B^-uB/B} \cap \overline{B^-vB/B}&\xrightarrow{\text{横截性}} \text{良定义的交点数} \\
&\xrightarrow{\text{Corollary 2.4}} \sum_w c^w_{u,v} \cdot [\overline{B^-wB/B}]_T \\
&\xrightarrow{\text{多项式代表元}} \mathfrak{S}_u \cdot \mathfrak{S}_v = \sum_w c^w_{u,v} \cdot \mathfrak{S}_w
\end{align}

\subsection{Graham 定理为何依赖几何性质}\label{sec:qa-graham-proof-geometry}

\textbf{问}：为什么 Graham 的证明要依赖几何性质？具体用了什么几何性质？到了量子上同调为何变得更复杂？

\textbf{答}：

\textbf{Graham 证明的核心几何思想}：
\begin{enumerate}
\item \textbf{几何相交}：在旗流形 $G/B$ 中，两个 Schubert 胞腔 $B^-uB/B$ 和 $B^-vB/B$ 的交集（在一般位置下）是若干 Schubert 胞腔的并。
\item \textbf{横截性保证非负}：由 Kleiman-Bertini 定理，当群作用足够一般时，相交是横截的（transverse），这保证了交集点的个数是良定义的非负整数。
\item \textbf{等变量子上同调的类}：在等变量子上同调中，Schubert 胞腔闭包 $\overline{B^-wB/B}$ 对应一个类 $[\overline{B^-wB/B}]_T \in H_T^*(G/B)$。
\item \textbf{类乘积 = 相交}：类乘法在几何上对应于胞腔的相交。
\item \textbf{多项式代表元}：等变 Schubert 类可以用标准 Schubert 多项式 $\mathfrak{S}_w(\mathbf{x})$ 作为代表元，这就将几何相交翻译为了多项式系数。
\end{enumerate}

\textbf{为何等变量子上同调比普通上同调更好用？}
\begin{itemize}
\item \textbf{整数系数丢失信息}：普通上同调只告诉我们"有多少个交点"，但不区分不同"类型"的交点。
\item \textbf{等变量子上同调保留权的信息}：等变类携带了 torus 作用的信息——让我们能看到系数中 $x_i$ 的具体结构，而不只是总数。
\end{itemize}

\textbf{为何量子上同调更复杂？}
\begin{enumerate}
\item \textbf{曲线代替点}：不再是计数点与点的交，而是计数满足特定 Gromov-Witten 不变量的曲线（稳定映射）。
\item \textbf{额外自由度}：稳定映射空间 $\overline{M}_{g,n}(G/P, d)$ 的维数依赖于次数 $d$，这引入了额外的形式变量 $q^d$。
\item \textbf{非横截情形出现}：在量子情形中，Schubert 胞腔的"相交"可能是整条曲线的重合。
\end{enumerate}

Mihalcea 的策略是：虽然量子上同调更复杂，但我们可以\textbf{退化}（degeneration）到经典情形。让某些参数趋于无穷大，量子上同调的乘法就退化回经典等变量子上同调的乘法。

\subsection{$B^-(w)$ 与 $B^-$ 的关系}\label{sec:qa-bminus-w-vs-bminus}

\textbf{问}：$B^-(w)$ 是什么？它和 $B^-$ 有什么关系？

\textbf{答}：

\textbf{定义回顾}：$$N^-(w) := N^- \cap w N^- w^{-1}$$$$B^-(w) := T \cdot N^-(w)$$

\textbf{各符号的含义}：
\begin{itemize}
\item $B^-$：相反 Borel 子群（下三角可逆矩阵群）
\item $N^-$：$B^-$ 的幂幺根基（单位下三角矩阵）
\item $T$：极大环面（对角矩阵群，$T = B \cap B^-$）
\item $N^-(w)$：$N^-$ 中"固定" $w$ 的子群
\item $B^-(w)$：由 $T$ 和 $N^-(w)$ 生成的子群（半直积）
\end{itemize}

\textbf{极端情形}：
\begin{itemize}
\item 当 $w = w_0$（最长元）时：$B^-(w_0) = T$（因为 $N^-(w_0) = \{e\}$）
\item 当 $w = e$（单位元）时：$B^-(e) = B^-$
\end{itemize}

\textbf{为什么 Refined Graham Positivity 更强}：原始 Graham 要求 $B^-$-不变闭链，而 Refined 版本要求 $B^-(w)$-不变闭链。条件更弱（更容易满足），但结论反而更强。

\subsection{$\mathbb{N}[-\alpha]_{\alpha \in I(w)}$ 的含义}\label{sec:qa-n-minus-alpha}

\textbf{问}：在 Refined Graham Positivity 的证明中，系数 $c^w_{u,v}(\mathbf{y}, \mathbf{t}) \in \mathbb{N}[-\alpha]_{\alpha \in I(\tau)}$ 这个符号是什么意思？为什么用 $-\alpha$ 而不是 $\alpha$？

\textbf{答}：

\textbf{符号分解}：记号 $\mathbb{N}[-\alpha]_{\alpha \in I(\tau)}$ 表示"以 $\{-\alpha \mid \alpha \in I(\tau)\}$ 为变量的非负整数系数多项式"。

\textbf{具体例子}：由计算，$I(\tau) = \{y_j - t_i \mid 1 \leq i,j \leq n\}$，因此$$-\alpha = -(y_j - t_i) = t_i - y_j$$所以 $\mathbb{N}[-\alpha]_{\alpha \in I(\tau)} = \mathbb{N}[t_i - y_j]_{i,j \geq 1}$，这正是 Graham Positivity 断言的系数形式。

\textbf{为什么是 $-\alpha$？}
\begin{enumerate}
\item \textbf{符号约定}：在等变量子上同调中，$T$ 的权通常用负根表示。具体来说，$H_T^*(\mathfrak{pt}) = \mathbb{Z}[t_1, \ldots, t_n]$ 中 $t_i$ 对应负根 $-\epsilon_i$。
\item \textbf{Corollary 2.4 的形式}：其中 $\mathbb{N}[-\alpha]$ 表示系数是关于负根的非负多项式。
\end{enumerate}

\textbf{直观理解}：这是把"关于 $t_i - y_j$ 的非负整数系数多项式"写得更加抽象和紧凑的形式。两种写法完全等价。

\subsection{EQLR 系数的定义与意义}\label{sec:qa-eqlr-coefficient}

\textbf{问}：EQLR 系数是什么？它在等变量量子上同调中扮演什么角色？

\textbf{答}：

\textbf{EQLR = Equivariant Quantum Littlewood-Richardson（等变量子 Littlewood-Richardson）}

EQLR 系数 $c_{u,v}^{w,d}$ 是等变量量子上同调 $QH_T^*(X)$ 中的结构常数，由以下展开式定义：$$\sigma(u)^T \circ \sigma(v)^T = \sum_{w,d} q^d \cdot c_{u,v}^{w,d} \cdot \sigma(w)^T$$其中：
\begin{itemize}
\item $\sigma(u)^T$ 是 $T$-等变 Schubert 类
\item $q^d$ 是量子参数（对应 Gromov-Witten 不变量）
\item $d$ 是多重次数（quantum degree）
\end{itemize}
